\chapter{Finding V --- attention versus memory, and the cue-attention question}
\label{ch:attention}

The deep dive raised a mechanistic question that is, at the time of writing, the program's
live frontier: in the canonical multiplicatively-gated model, the cueing \emph{behaviour}
is real and graded, yet the per-frame attention map does not orient to the cue. This
chapter sets out the dissociation, the mechanism behind it, and the matched family of
architectures we are training to resolve it.

\section{The cue modulates behaviour}
First, to be unambiguous: the cue strongly modulates the model's behaviour. The detection
threshold for a cued change is up to $7.5^\circ$ lower than for an uncued change, and the
benefit scales with validity (Chapter~\ref{ch:deepdive}). The model reads the cue, holds
it, and uses it. The question is not \emph{whether} the cue is used but \emph{through what
internal variable}.

\section{The attention is change-locked, not cue-locked}
We measured the four-patch attention map at every frame for a $100\%$ cue at $S_1$ with a
$25^\circ$ change at $S_1$ (Fig.~\ref{fig:attn_change}). The attention is uniform
($\approx0.25$ per patch) through the cue frame and the entire delay, and it sharpens onto
$S_1$ only at frame $5$ --- the moment the change appears. The full $4\times4$ query-key
matrix at the cue frame is essentially uniform, so this is not an artefact of how the map
is summarised: there is no cue-driven structure to summarise. A no-change control confirms
the reading (Fig.~\ref{fig:attn_nochange}): with the same $100\%$ cue but no change, the
attention stays flat across all seven frames (maximum deviation from uniform $0.039$). The
attention fires at a location when, and only when, that location \emph{changes}.

\begin{figure}[t]
\centering
\includegraphics[width=0.95\linewidth]{fig_attention_maps_S1.pdf}
\caption{Per-frame attention (top: input frames; bottom: $2\times2$ attention map, S1 =
top-left) for a $100\%$ cue at $S_1$ with a $25^\circ$ change at $S_1$. The cue (frame 1)
draws no attention; the attention concentrates on $S_1$ only at the change (frame 5).}
\label{fig:attn_change}
\end{figure}

\begin{figure}[t]
\centering
\includegraphics[width=0.85\linewidth]{fig_attention_maps_S1_nochange.png}
\caption{No-change control, $500$ trials, $100\%$ cue at $S_1$. With nothing to lock onto
the attention stays uniform across every frame --- the cue alone never moves it.}
\label{fig:attn_nochange}
\end{figure}

\section{Where the cue actually lives}
If not in the attention, where? In the recurrent memory. Decoding from the flattened memory
$\mathbf H$ at the pre-change frame, both the cued \emph{location} ($S_1$ vs.\ $S_4$) and
the \emph{validity} ($25/50/75/100\%$) are recovered at $100\%$ accuracy (chance $50\%$ and
$25\%$). The cue and its reliability are perfectly maintained in working memory, where they
set the decision threshold; the attention, meanwhile, is a bottom-up change detector. This
is a clean dissociation: \textbf{cue $\to$ behaviour: yes; cue $\to$ attention map: no}.

\section{The mechanism}
The dissociation has a transparent cause in the multiplicative form
(Eq.~\ref{eq:mult}). Because $\mathbf Q=(\X\mathbf W_{XQ})\odot(\Hone\mathbf W_{HQ})$, the
query magnitude is gated by the memory projection. At the cue frame the memory has not yet
built up, the gate $\lvert\Hone\mathbf W_{HQ}\rvert$ is near zero ($\approx0.02$), and the
queries and keys collapse toward zero, so the attention is uniform. The model is then
\emph{free} to keep that gate small, because the cue reaches memory through the image
residual $\X$ in $\mathbf Z=\X+\mathbf V_{\text{filt}}$ and the LSTM, without the attention
ever needing to orient. The multiplicative form, in short, \emph{admits} a memory-only
solution, and on the seven-step task with a frozen VAE and a memory readout, the harness
finds it.

\section{The open question}
This is where honesty is required. The lab's other models --- multi-head, cross-attention
variants in the spirit of \code{v5} and \code{v11} --- reliably \emph{do} orient attention
to the cue, across many trained networks. So orienting is clearly achievable with these
mechanisms; the multiplicative-gated reproduction is the outlier. The cause is therefore
\emph{some feature of our specific reproduction} --- the gate form, the head count, the
memory readout, the harness, the frozen front-end --- not a property of memory-gated
attention in general. We do not yet know which feature, and we are not guessing: we are
testing it directly with a matched family of architectures.

\section{The matched family under test}
Each variant changes exactly one thing about the attention/memory wiring on the identical
base (frozen VAE, harness, readout dimension), so any difference in whether attention
orients to the cue is attributable to that change. Table~\ref{tab:family} lists them.

\begin{table}[t]
\centering
\caption{The matched family of attention/memory wirings under training to resolve the
cue-attention question. All share the frozen-VAE front-end, the PAC$+$QR-DQN$+$PER harness,
and the $4096$-d readout.}
\label{tab:family}
\small
\begin{tabular}{p{2.2cm}p{7.2cm}p{2.0cm}}
\toprule
Variant & What it changes about the feedback / memory & Status \\
\midrule
multiplicative & canonical $\mathbf Q=(\X\mathbf W)\odot(\Hone\mathbf W)$ (the reproduction) & works; change-locked \\
affine & memory derives a scale-\emph{matrix} $+$ shift: $\X'=\Gamma(\Hone)\X+\beta(\Hone)$ & \textbf{failed (always-wait)} \\
dualhead & a memory-free sensory head $\parallel$ the multiplicative feedback head & works; weak orient \\
two-lstm & split memory: $\Hone$ feeds attention, $\Htwo$($\Hone$) feeds heads & \textbf{best detector} \\
crossattn & additive cross-talk: $\mathbf K,\mathbf V=[\X\parallel\Hone]$, two LSTMs from $\mathbf Z$ & works; memory-routed \\
dualmem & memories \emph{query}: $\mathbf Q([\Hone\parallel\Htwo])$, $\mathbf K,\mathbf V=[\X\parallel\Hone\parallel\Htwo]$ & works; memory-routed \\
\bottomrule
\end{tabular}
\end{table}

\section{Results: the family resolves the question}
All six variants are now trained to $\sim$$18{,}750$ iterations on the identical frozen-VAE
base, and each was evaluated greedily for behaviour ($250$ episodes on \code{validity4}) and
for cue-orienting with the same $100\%$-cue protocol used above: a fully reliable cue, the
attention to the \emph{cued} patch recorded across all seven frames, in a no-change
condition (so any pre-change bias toward the cued location is cue-driven, not change-driven)
and a change condition (so the change-locked component is visible). The
\emph{cue-orienting index} is the cued-patch attention at the pre-change stimulus frames
($t_{3,4}$) minus uniform; the \emph{change-locked spike} is the cued-patch attention jump
at the change frame ($t_5$). Table~\ref{tab:familyresults} reports both alongside behaviour;
Figures~\ref{fig:varattn} and~\ref{fig:varbeh} plot them.

\begin{table}[t]
\centering
\caption{Matched-family results. Behaviour is greedy on \code{validity4} ($T=7$);
cue-orient $>0$ means the attention is biased toward the cued location \emph{before} the
change (cue-driven). The headline: no variant produces a large positive cue-orient, while
two of them (\code{crossattn}, \code{dualmem}) attend markedly \emph{away} from the cued
image patch --- the cueing effect is carried by memory across the whole family.}
\label{tab:familyresults}
\small
\begin{tabular}{lccccrr}
\toprule
Variant & correct & hit & FA & median RT & cue-orient & chg-spike \\
\midrule
multiplicative & $0.804$ & $0.705$ & $0.091$ & $5.0$ & $-0.011$ & $+0.143$ \\
affine         & $0.504$ & $0.000$ & $0.000$ & ---   & $+0.014$ & $-0.001$ \\
dualhead       & $0.820$ & $0.708$ & $0.088$ & $5.0$ & $+0.038$ & $-0.002$ \\
two-lstm       & $\mathbf{0.872}$ & $\mathbf{0.836}$ & $0.094$ & $5.0$ & $+0.008$ & $+0.216$ \\
crossattn      & $0.840$ & $0.760$ & $0.080$ & $5.0$ & $-0.179$ & $-0.011$ \\
dualmem        & $0.828$ & $0.750$ & $0.079$ & $5.0$ & $-0.174$ & $+0.002$ \\
\bottomrule
\end{tabular}
\end{table}

\begin{figure}[t]
\centering
\includegraphics[width=0.99\linewidth]{fig_variants_attention.pdf}
\caption{\emph{Left:} attention to the cued patch across the seven frames, no-change,
$100\%$ cue, for every working variant. All hover near or below uniform ($0.25$) through the
cue and delay; none lifts the cued location before the change. \emph{Right:} the
cue-orienting index ($t_{3,4}$, solid) and the change-locked spike ($t_5$, hatched). The
models that read a direct image$\to$attention path (\code{multiplicative}, \code{two-lstm})
show a clear change-lock and no orienting; \code{crossattn} and \code{dualmem} sit well
below uniform on the cued image patch --- they route the cued location through memory rather
than attending it.}
\label{fig:varattn}
\end{figure}

\begin{figure}[t]
\centering
\includegraphics[width=0.62\linewidth]{fig_variants_behavior.pdf}
\caption{Behaviour on \code{validity4}. Five of six variants learn the task; \code{affine}
collapses to always-wait (hit $=0$, no false alarms, $P(\text{press})\le0.004$ even on a
clear $60^\circ$ change). \code{two-lstm} is the best detector.}
\label{fig:varbeh}
\end{figure}

The matched family answers the question, and not in the direction the single early data
point suggested. \textbf{No variant restores strong cue-orienting attention on this task.}
The largest positive cue-orienting index is \code{dualhead}'s sensory head at $+0.038$ ---
real but small, and too weak to carry an attention-allocation figure. The
\code{multiplicative} and \code{two-lstm} models are change-locked exactly as the deep dive
described (spike $+0.143$, $+0.216$; orient $\approx0$). And the two richest cross-talk
wirings, \code{crossattn} and \code{dualmem}, attend \emph{below} uniform at the cued image
location ($-0.179$, $-0.174$): given the choice between image keys and memory keys, they
spend their attention on memory, leaving the bottom-up image patch under-weighted. The
dissociation of Chapter~\ref{ch:attention} is therefore not a quirk of the multiplicative
gate; it is \emph{general} across six attention/memory wirings on this perception front-end
and this task. Reward, on the seven-step task, does not make attention orient to the cue ---
it makes \emph{memory} hold the cue.

\paragraph{A correction.} An earlier reading of a single partially-trained \code{affine}
checkpoint suggested it ``already attends the cue.'' That was premature: fully evaluated,
\code{affine} never presses (Fig.~\ref{fig:varbeh}), so its attention is not behaviourally
meaningful, and the apparent cue structure does not survive into a model that performs the
task. The lesson --- verify behaviour before reading an attention map --- is the same one
the dissociation itself teaches.

\section{Why this matters}
For the empirical manuscript, the behavioural reproduction (Chapter~\ref{ch:deepdive})
stands regardless. But the result reshapes the attention-map story in two useful ways.
First, it converts the open question into a \emph{finding}: across this family, the cueing
effect is memory-borne, and the attention map is a bottom-up change detector --- a
defensible, falsifiable claim now backed by six wirings rather than one. Second, it picks
the network: \code{two-lstm} is both the best detector and a minimal, paper-faithful change
to the published multiplicative attention (a second stacked memory), so it is the natural
network to carry forward, with the memory-driven account of cueing stated plainly rather
than hidden behind an attention figure that the mechanism does not support. If a future
task with a longer delay or a harder cue \emph{forces} pre-orienting, this same family is
the instrument to re-run; on the seven-step task, the honest figure is the change-locked
one.
